\clearpage\fancyhead[L]{\small ANZEIGEN. LOGARITHMISCHE TAFELN.}
\SourceLeaf{262}
\pdfbookmark[1]{Babbage. 1828 Januar 19.}{babbage1828}\par\Needspace{10\baselineskip}\begin{center}\rule{0.7\linewidth}{0.4pt}\par Göttingische gelehrte Anzeigen. 1828 Januar 19.\par\rule{0.7\linewidth}{0.4pt}\end{center}

\noindent\textit{Table of logarithms of the natural numbers, from} 1 \textit{to} 108000, \textit{by} \textsc{Charles}
\textsc{Babbage.} \textit{Stereotyped. London. Printed for J. Mawman.} 1827. 202 S. gr. 8.

Dieser neue Abdruck der Logarithmentafeln zeichnet sich vor andern durch
eine geflissentlichere Beachtung kleiner Nebenumstände aus. Wer nur von Zeit
zu Zeit einmal veranlasst wird, einige Logarithmen in den Tafeln aufzusuchen,
verlangt von ihnen hauptsächlich nur möglich grösste Correctheit. Allein für andere, denen die Tafeln ein tägliches Arbeitsgeräth sind, bleiben auch die geringfügigsten Umstände, die auf die Bequemlichkeit des Gebrauchs Einfluss haben
können, nicht mehr gleichgültig. Farbe, Stärke und Schönheit des Papiers;
Format; Grösse, Schärfe und gefälliger Schnitt der Typen: Beschaffenheit der
Druckerschwärze; Anordnung der Zahlen, um das was man sucht ohne Ermüdung des Auges schnell und sicher zu finden; Vorhandensein von allem, was man
braucht, aber auch Abwesenheit von allem, was man nicht brauchen mag, und
was sonst die leichte Uebersicht nur stören würde, alle diese Umstände erhalten
eine gewisse Wichtigkeit bei einem Geschäfte, welches man täglich hundert mal
wiederholt. Freilich hängt dabei manches von der Individualität des Rechnenden
und von seiner Gewohnheit ab, so dass nicht wohl Eine Ausgabe allen am besten
gefallen kann; dem Kurzsichtigen ist ein grosses Format beschwerlich, und er
zieht kleinere Typen vor, während es sich bei dem Weitsichtigen umgekehrt verhält; der weniger geübte Rechner legt einen Werth auf diesen oder jenen Zusatz,
welchen der geübtere, der mehrere Theile der Operationen ohne Anstrengung im
Kopf macht, lieber wegwünscht, weil alles Ueberflüssige nur störend wirkt. Inzwischen gibt es doch auch allgemeingültige Regeln. Der Herausgeber der vorliegenden Tafeln hat in der Vorrede ein Dutzend solcher Vorschriften zusammengestellt, die er aus der Vergleichung vieler Logarithmentafeln abgeleitet hat, und
die meistens sogleich von selbst einleuchten, obwohl einige davon nur unter Einschränkungen anzuerkennen sein möchten.

Was die gegenwärtige Ausgabe der Logarithmentafeln am meisten von andern unterscheidet, ist, dass sie auf farbiges (gelbes) Papier abgedruckt ist: man
gewöhnt sich daran bald, und findet es, besonders zum Gebrauche bei Licht, an%
\SourceLeaf{263}genehmer als weisses. Die Anordnung ist im Wesentlichen die gewöhnliche.
Nach einem Wechsel der dritten Ziffer ist für die übrigen Logarithmen in derselben Zeile die vierte Ziffer mit andern Typen, etwa halb so gross wie die übrigen,
gesetzt. Für Ref., der überhaupt auf solche Warnungszeichen wenig Werth setzt,
haben diese Typen, auch nach einem Gebrauch von ein paar Monaten, noch nicht
das fremdartig Störende verloren, und er würde den in \textsc{Vega}’s Tafel gebrauchten
Sternchen, oder den Punkten, die in die bessern Exemplare von \textsc{Taylor}’s Tafeln
eingedruckt sind, den Vorzug geben. Sternchen hat der Herausgeber desswegen
nicht gebrauchen wollen, weil sonst die Columnen zu breit geworden, und dann
für die Verwandlung der Zahlen, als Secunden betrachtet, in Grade, Minuten
und Secunden kein Raum geblieben wäre. Referent würde diese Verwandlungscolumne (eben so eingerichtet, wie man sie aus \textsc{Callet}’s Tafeln kennt) gern entbehrt haben; ein geübter Rechner wird nicht, einer so leicht selbst im Kopfe zu
machenden Verwandlung wegen, erst die Tafeln aufblättern; zweckmässig ist es
aber, dass der Herausgeber diese Columnen wenigstens durch eine starke Linie
von der Logarithmentafel geschieden hat. Endlich ist der Fall wo die letzte Ziffer eine Vergrösserung erlitten hat, (weil der weiter fortlaufende Logarithm als
achte Ziffer 5 oder eine grössere gehabt haben würde) durch einen unter diese siebente Ziffer gesetzten Punkt ausgezeichnet; diese Einrichtung ist wenigstens besser, als die von \textsc{Prasse} gewählte, durch Typen von anderer Form; doch behalten
auch so einige Ziffern noch etwas unangenehm Fremdartiges, was nach unserer
schon bei Anzeige der \textsc{Prasse}’schen Tafeln in diesen Blättern [1811 Mai 25] erwähnten Ansicht durch den Nutzen nicht aufgewogen wird.

Uebrigens lassen Typen und Papier bei dieser Ausgabe der Logarithmentafeln nichts zu wünschen übrig, und auf die Correctheit ist die ausgezeichnetste
Sorgfalt verwandt. Bei einer Ausgabe, wo auf die kleinen die Bequemlichkeit
des Gebrauchs angehenden Umstände so viel Sorgfalt verwandt ist, fällt die spärliche Ansetzung der Proportionaltheile auf; häufig muss man, um diejenigen zu
finden, die man eben braucht, voraus- oder zurückblättern.
\SourceLeaf{264}
\pdfbookmark[1]{Hassler. 1831 März 31.}{hassler1831}\par\Needspace{10\baselineskip}\begin{center}\rule{0.7\linewidth}{0.4pt}\par Göttingische gelehrte Anzeigen. 1831 März 31.\par\rule{0.7\linewidth}{0.4pt}\end{center}

\noindent\textit{Tabulae logarithmicae et trigonometricae notis septem decimalibus expressae.}
\textit{In forma minima. Purgatae ab erroribus praecedentium tabularum cura} F. K. \textsc{Hassler.} New York. G. C. und H. Carvill.

Diese transatlantische Ausgabe der Logarithmentafeln zeichnet sich durch
eine ganz vorzügliche Nettigkeit des stereotypisch ausgeführten Drucks aus. Sie
enthält, durchgehends auf sieben Decimalstellen, die Logarithmen der Zahlen
bis 100000, die Logarithmen der Sinus, Tangenten, Cosinus und Cotangenten
im ersten Grade durch alle Secunden, in den beiden folgenden von zehn zu zehn
Secunden, und für alle übrigen Grade des Quadranten von dreissig zu dreissig
Secunden. Ausserdem die natürlichen Sinus und Tangenten durchgehends von
dreissig zu dreissig Secunden. Und diess alles in dem Raum von 312 Seiten in
klein Octav oder gross Duodez. Eine solche Zusammendrängung war freilich nur
durch sehr kleine Typen zu erreichen, welche, bei aller Schönheit, doch wohl
für die meisten nicht kurzsichtigen Augen zum täglichen Gebrauch fast zu klein
sein möchten. Auf die Correctheit scheint eine ganz besondere Sorgfalt gewandt
zu sein, wenigstens ist uns bei dem eine Zeitlang versuchten häufigen Gebrauch
gar kein Druckfehler aufgestossen.

\pdfbookmark[1]{Auflösung quadratischer Gleichungen. 1840.}{quadratisch1840}\par\Needspace{10\baselineskip}\begin{center}\rule{0.7\linewidth}{0.4pt}\par \textsc{Vega} und \textsc{Hülsse.} Sammlung mathematischer Tafeln. 1840.\par\rule{0.7\linewidth}{0.4pt}\end{center}

\noindent\textit{Auflösung quadratischer Gleichungen in der Form, dass nicht die Coëfficienten}
\textit{der Gleichung selbst, sondern deren Logarithmen gegeben sind, und dass man auch}
\textit{nicht ihre Wurzeln selbst (oder eine derselben), sondern vielmehr deren Logarithmen}
\textit{zu anderweitiger Benutzung nöthig hat.}

Die Ausführung dieses Geschäfts, bloss mit Hülfe der gewöhnlichen Logarithmentafeln, erfordert nothwendig ein vierfaches Aufschlagen in denselben; man
reicht mit einem zweifachen Aufschlagen aus, wenn man entweder die trigonometrischen Logarithmentafeln oder die Tafeln für Logarithmen von Summen und
\SourceLeaf{265}
Differenzen benutzt; man reicht mit einem einfachen Aufschlagen aus, wenn letztere Tafel, wie hier, durch die Zusatzcolumnen $D$, $E$ und $F$ erweitert ist.

Die vierte Columne, $D$, enthält $B+C$, die fünfte, $E$, enthält $A+C$,
die sechste, $F$, enthält $B-A$. Nehmen wir an, dass die Zahlen selbst, denen
die Logarithmen $A$, $B$, $C$ zugehören, beziehungsweise $a$, $b$, $c$ sind, so enthielte
also die vierte Columne die Logarithmen von $bc$, die fünfte die von $ac$, die
sechste die von $\frac{b}{a}$.

Es mag noch darauf aufmerksam gemacht werden, dass die Logarithmen in
der sechsten Columne $F$, welche bis zu $A=0{,}208$ positiv sind, von $A=0{,}209$
an negativ werden (S.\ 640); es ist ziemlich gleichgültig, ob sie in dieser negativen Form oder, wie hier, durch ihre Complemente angesetzt werden; also könnte
z.\ B.\ für $A=0{,}367$ unter $F$ stehen $-0{,}21180$ oder, wie hier, $9{,}78820$.

Für die Anwendung dieser Zusatzcolumnen auf die Auflösung der quadratischen Gleichung
\[
pxx+qx+r=0
\]
selbst müssen vier verschiedene Fälle unterschieden werden, nämlich
\begin{enumerate}
\item[I.] $p$ und $r$ haben gleiche Zeichen und $\frac{qq}{pr}$ ist nicht kleiner als $4$;
\item[II.] $p$ und $r$ haben gleiche Zeichen und $\frac{qq}{pr}$ ist kleiner als $4$;
\item[III.] $p$ und $r$ haben entgegengesetzte Zeichen und $-\frac{pr}{qq}$ ist grösser als $2$;
\item[IV.] $p$ und $r$ haben entgegengesetzte Zeichen und $-\frac{pr}{qq}$ ist kleiner als $2$.
\end{enumerate}
Der Fall, wo $-\frac{pr}{qq}=2$ ist, kann sowohl zu III.\ als zu IV.\ gezählt werden.

Im Falle II.\ sind die Wurzeln imaginär, in den übrigen erhält man jede der
beiden Wurzeln auf eine doppelte Art. Durch folgendes Schema ist alles leicht
zu übersehen, wobei
\[
\frac{q}{p}=h,\quad \frac{r}{q}=g
\]
gesetzt ist, theils zur Abkürzung, theils weil die Rechnung wirklich in dieser
Form am bequemsten geführt wird.

\begin{center}
\begingroup\renewcommand{\arraystretch}{1.7}\begin{tabular}{rlll}
 & & \emph{Erste Wurzel.} & \emph{Zweite Wurzel.} \\
I. & $+\frac{h}{g}=bc=d$ & $-\frac{h}{b}=-gc$ & $-gb=-\frac{h}{c}$ \\
II. & $+\frac{h}{g}<4$ & Imaginär. & Imaginär. \\
III. & $-\frac{g}{h}=ac=e$ & $+ha=-\frac{g}{c}$ & $+\frac{g}{a}=-hc$ \\
IV. & $-\frac{g}{h}=\frac{b}{a}=f$ & $+\frac{h}{a}=-\frac{g}{b}$ & $+ga=-hb$
\end{tabular}\endgroup
\end{center}
\par
\SourceLeaf{266}
Die Beweise der Vorschriften wird sich jeder leicht selbst entwickeln können. Ein Beispiel des Gebrauchs mag zum Ueberfluss noch hergesetzt werden.

\begin{center}
\begin{tabular}{rl}
Es sei gegeben & $\log p=0{,}69897$ \\
 & $\log q=0{,}84510$ \\
 & $\log r=0{,}77815\,\mathrm{n}$ \\
Also & $\log h=0{,}14613$ \\
 & $\log g=9{,}93305\,\mathrm{n}$ \\
 & $\log(-\frac{g}{h})=9{,}78692$
\end{tabular}
\end{center}

Man sieht sogleich, dass hier der Fall IV.\ statt findet und also $9{,}78692$ in
der sechsten Columne unter $B-A$ oder $F$ zu suchen sein wird. Was von der
Tafel (S.\ 643) hier nöthig wäre ist nur Folgendes:
\begin{center}
\begin{tabular}{c|c|c}
$A$ & $B$ & $B-A$ oder $F$ \\
$0{,}367$ & $0{,}15520$ & $9{,}78820$ \\
$0{,}368$ & $0{,}15489$ & $9{,}78689$
\end{tabular}
\end{center}
Zu $F=9{,}78692$ gehört also
\begin{center}
\begin{tabular}{l}
$A=0{,}36798=\log a$, woraus $\log\frac{h}{a}=9{,}77815$, $\log ga=0{,}30103\,\mathrm{n}$ oder \\
$B=0{,}15490=\log b$, woraus $\log(-\frac{g}{b})=9{,}77815$, $\log(-hb)=0{,}30103\,\mathrm{n}$.
\end{tabular}
\end{center}

\par\Needspace{12\baselineskip}\pdfbookmark[1]{Vega. Thesaurus Logarithmorum. 1851.}{vega1851}
\begin{center}
\rule{0.7\linewidth}{0.4pt}\par
Astronomische Nachrichten. Nr.\ 756. 1851. Mai 2.\par
\rule{0.7\linewidth}{0.4pt}
\end{center}
\begin{center}\emph{Einige Bemerkungen zu} \textsc{Vega}'s \emph{Thesaurus Logarithmorum, von Herrn}\par
\emph{Geheimen Hofrath} \textsc{Gauss}.\end{center}

Der \emph{Thesaurus Logarithmorum} von \textsc{Vega} ist bekanntlich seinem grössten
Theile nach ein neuer Abdruck der grössern \textsc{Vlacq}'schen Logarithmentafeln. In
der Vorrede führt \textsc{Vega} eine nicht unbeträchtliche Anzahl von Fehlern im Original an, die er verbessert hat, mit dem Zusatz, dass er ausser diesen noch eine
sehr grosse Menge von Unrichtigkeiten an der letzten Ziffer der Logarithmen berichtigt habe, zu dem Betrage von einer, zwei, drei, vier Einheiten. Mit gleicher Sorgfalt seien auch die neu hinzugekommenen Tafeln (namentlich also die in
\SourceLeaf{267}
den beiden ersten Graden für alle einzelnen Secunden angegebenen Logarithmen
der trigonometrischen Grössen) berechnet, geprüft und berichtigt. \textsc{Vega} scheint
nun mit der Hoffnung sich geschmeichelt zu haben, dass auf diese Weise seine
Tafeln fast fehlerfrei geworden seien, und verspricht, um zu \emph{vollkommen} fehlerfreien Tafeln zu gelangen, für die erste Anzeige jedes etwa noch stehen gebliebenen Fehlers, der zu falscher Rechnung Anlass geben könne (\emph{pro sphalmatibus calculum turbantibus}) eine Prämie von einem Ducaten zu bezahlen. Ob diese vom
$1^{\mathrm{sten}}$ October $1794$ datirte Ausgelobung jemals Folge gehabt hat, ist mir nicht
bekannt.

Es ist mir zweifelhaft, ob \textsc{Vega} sich ganz klar gemacht habe, was für Fehler als möglicher Weise zu falschen Rechnungen Anlass gebend betrachtet werden sollten. Für alle Tafeln, welche bestimmt sind, theoretisch feststehende
irrationale Grössen darzustellen, gilt bekanntlich der Grundsatz, dass die Tabulargrösse dem wahren Werthe allemal so nahe kommen soll, als bei der gewählten Anzahl von Decimalstellen möglich ist, und es darf folglich die Abweichung
niemals mehr als eine halbe Einheit der letzten Decimale betragen. Jeder Verstoss gegen diese strenge Norm ist ein Fehler, der möglicherweise einen sich auf
die strenge Uebereinstimmung verlassenden Rechner zu einem unrichtigen Resultate verleiten kann. Lässt man diese strenge Auslegung fahren, und mischt in
sein Urtheil eine Rücksicht auf \emph{Erheblichkeit} der Unrichtigkeit ein, so verirrt
sich die Entscheidung in das Gebiet der Willkühr. Der schon vorhin erwähnte
Umstand, dass \textsc{Vega} selbst von Correctionen an den \textsc{Vlacq}'schen Tafeln spricht,
die nur eine Einheit in der letzten Stelle betrugen, und dass er vollkommene Fehlerfreiheit wie sein Ziel bezeichnet, scheint allerdings darauf hinzudeuten, dass
er die strenge Auslegung im Sinne gehabt. Auch habe ich den ersten Theil des
Thesaurus, der die Logarithmen der laufenden ganzen Zahlen enthält, bei sehr
vielen gelegentlich gemachten Vergleichungen mit mehrstelligen Bestimmungen
immer sehr correct gefunden.

Es sind seit jener Zeit bei mehrern andern logarithmischen Tafeln, in der
Absicht ihre Correctheit zu vervollständigen, ähnliche Ausbietungen von Preisen
für die erste Anzeige von Fehlern in den Zahlen gemacht: ich weiss jedoch nicht,
ob dieselben einen Erfolg gemacht haben, mit Ausnahme des bei Tauchnitz in
Leipzig $1847$ von \textsc{Köhler} herausgegebenen logarithmisch-trigonometrischen Handbuchs. Der Verleger dieser Tafeln versprach bei dem ersten Erscheinen, für die erste
\SourceLeaf{268}
binnen einer gesetzten Frist eingesandte Anzeige eines jeden Fehlers, welcher
falsche Resultate veranlassen könne, einen Louisd'or zu bezahlen, und nach einem gedruckten Bericht vom 1. Juli 1848 ist diese Prämie für vier zur Anzeige
gebrachte Fehler wirklich ausgezahlt. Was nun dabei eine ehrende Erwähnung
verdient, ist der Umstand, dass dem einen dieser Fehler jene Qualification nur
unter Anerkennung obiger strengen Auslegung zugesprochen werden konnte. Es
war nemlich der Logarithm von \(103000\), welcher, auf \(12\) Stellen genau,
\[
=5{,}012837224705
\]
ist, in den ersten Abdrücken mit acht Ziffern
\[
=5{,}01283723
\]
angesetzt, während die principmässig abgekürzte Zahl \(5{,}01283722\) ist. Es ist
zu wünschen, dass in künftig vorkommenden ähnlichen Fällen, diese Entscheidung als Präcedenz respectirt werde.

Bekannte sich \textsc{Vega} zu derselben strengen Auslegung. so hätte es ihm leicht
gehen mögen, wie dem König Shiram, dessen Kornkammern nicht ausreichten,
dem Erfinder des Schachspiels die ihm zugesagte Belohnung zu gewähren.

Dass die von \textsc{Vega} ausgebotene Belohnung, unter jener Voraussetzung. ihm
theuer zu stehen kommen konnte, lässt sich schon, ohne alles Nachrechnen, aus
einem Umstande erkennen, der leicht zu bestätigen, jedoch meines Wissens anderweit noch nirgends zur Sprache gebracht ist. Dieser Umstand besteht darin,
dass in der Tafel für die Logarithmen der trigonometrischen Grössen die Zahlen
der Sinuscolumne fast ohne Ausnahme\HistoricalNote{21} der Summe der Zahlen der Cosinuscolumne und der Tangentencolumne genau gleich sind. Da alle diese Zahlen nur
abgekürzte Werthe der irrationalen genauen Grössen sind, so ist klar, dass bei
streng richtiger Abkürzung jene Gleichheit nicht Statt finden wird, in allen den
Fällen, wo die Abweichungen von den genauen Werthen in der zweiten und dritten Columne gleiche Zeichen haben, und ihre Summe mehr als eine halbe Einheit der letzten Decimale beträgt. Man übersieht leicht, dass bei einer grossen
\SourceLeaf{269}
Menge von Fällen dieser Ausnahmefall durchschnittlich einmal unter vieren vorkommen wird, was sich bei wirklicher Abzählung in solchen siebenziffrigen Tafeln, wo auf die Richtigkeit der letzten Ziffer mit Sorgfalt gehalten ist, bestätigt
findet. (Es hat z. B. eine solche Abzählung in den erwähnten \textsc{Köhler}'schen Tafeln an den \(900\) einzelnen Minuten von \(30^{0}0'\) bis \(44^{0}59'\) genau \(225\) solcher Fälle
ergeben, welche vollkommene Uebereinstimmung allerdings für zufällig zu halten ist). \textsc{Vega}'s Thesaurus enthält die Logarithmen der Sinus, Cosinus und Tangenten von \(22680\) Bogen. Unter der Voraussetzung also, dass alle Zahlen in
zweien dieser Columnen scharf nach dem Princip abgekürzt seien, wird man, mit
einer geringen Unsicherheit im Mehr oder Weniger, in der dritten Columne \(5670\)
Logarithmen erwarten dürfen, die um eine Einheit unrichtig angesetzt sind. Für
den Preis von so vielen Ducaten würden wohl befähigte Personen zur vollständigen Neurechnung bereit gewesen sein.

Welche zwei Columnen die ursprünglichen sind, wird man aus der Vergleichung mit anderweitigen auf mehr als zehn Stellen zuverlässigen Bestimmungen erkennen können, schon an wenigen Bögen und mit Gewissheit, wenn die
Voraussetzung der strengen Richtigkeit der ursprünglichen Columnen zutrifft, im
entgegengesetzten Fall aber an einer etwas beträchtlichern Menge wenigstens mit
überwiegender Wahrscheinlichkeit.

Einiges hiezu dienliche findet man schon fertig vor am Schluss der Decimaltafeln, welche \textsc{Hobert} und \textsc{Ideler} 1799 geliefert haben, und die zwar nur mit
sieben Ziffern abgedruckt sind, aber in der Handschrift auf doppelt so viele vollendet waren. Es werden daselbst \(138\) fehlerhafte Logarithmen des \textsc{Vega}'schen
Thesaurus angezeigt, von denen \(127\), um genau gesetzmässig zu werden, einer
Correction von einer Einheit in der zehnten Stelle bedürfen, \(10\) einer Correction von
zwei Einheiten, und eine drei Einheiten. Von den \(127\) fallen \(49\) auf Sinus, \(35\)
auf Cosinus und \(43\) auf Tangenten; die übrigen \(11\) Correctionen, von \(2\) oder \(3\)
Einheiten, beziehen sich blos auf Tangenten. Die betreffenden Bögen sind die
mit \(27\) Minuten messbaren, und die nicht mit angeführten sind diejenigen, wo
die Logarithmen keiner Correctionen bedurften. Es hätten übrigens diese Vergleichungen auch ohne die \textsc{Hobert-Ideler}'schen handschriftlichen Tafeln gemacht
werden können, da die allgemein verbreiteten \textsc{Callet}'schen Tafeln die Logarithmen der Sinus und Cosinus auf \(14\) Ziffern für alle Tausendtheile des Quadranten
enthalten. Wenn man aus diesen auch noch das Nöthige für alle Bögen unter
\SourceLeaf{270}
zwei Grad entlehnt, die durch \(5'24''\) aber nicht durch \(27'\) messbar sind, so gewinnt man die Vergleichung \textsc{Vega}'scher Logarithmen von \(118\) Bögen mit schärfern Bestimmungen, wovon die Resultate in folgendem Abriss (I) zusammengestellt sind:
\[
\begin{array}{rccc}
 & \mathrm{Sin.} & \mathrm{Cos.} & \mathrm{Tang.} \\
0 & 65 & 75 & 54 \\
1 & 53 & 43 & 53 \\
2 &    &    & 10 \\
3 &    &    & 1
\end{array}
\]
Die Bedeutung dieser Zahlen, z. B. der in der letzten Columne ist, dass unter
\(118\) Tangenten-Logarithmen \(54\) richtig angesetzt sind, \(53\) einer Correction von
einer Einheit in der zehnten Stelle bedürfen, \(10\) einer Correction von zwei Einheiten, und einer der Correction von drei Einheiten. Es folgt hieraus schon entschieden, dass \emph{keine} Columne der \textsc{Vega}'schen Tafel durchaus richtig angesetzt
ist, und mit überwiegender Wahrscheinlichkeit, dass die Tangenten-Logarithmen
die abgeleiteten sind, durch die einfache Subtraction der Zahlen der Cosinuscolumne von denen der Sinuscolumne.

Die Verfasser der erwähnten Decimaltafeln hätten übrigens doppelt so viele
Vergleichungen aus ihrer Handschrift geben können. Man kann sich aber eine
noch viel grössere Ausbeute verschaffen, wenn man die höchst schätzbare Tafel
von \textsc{Brigg}'s (in der nach dessen Tode von \textsc{Gellibrand} 1633 herausgegebenen Trigonometria Britannica) benutzt, welche die Logarithmen der Sinus und Cosinus
für alle Hunderttheile der gewöhnlichen Grade auf \(14\) Ziffern liefert, und aus
welcher \textsc{Callet} die oben erwähnten Zahlen entlehnt hat. Sie enthält das Material, um \textsc{Vega}'s Tafeln bei \(1060\) Bögen, also zusammen \(3180\) Logarithmen, prüfen zu können. Ich selbst habe mich jedoch darauf beschränkt, diese Prüfung
an \(81\) Bögen oder an \(243\) Logarithmen, von \(14^{0}0'\) bis \(18^{0}0'\) vorzunehmen, wovon
das Resultat hier folgt (II):
\[
\begin{array}{rccc}
 & \mathrm{Sin.} & \mathrm{Cos.} & \mathrm{Tang.} \\
0 & 29 & 56 & 36 \\
1 & 52 & 25 & 42 \\
2 &    &    & 3
\end{array}
\]
\par
\SourceLeaf{271}
Wirft man die Gruppen (I) und (II) zusammen, so jedoch, dass man die
beiden gemeinschaftlichen Bögen nur einmal in Rechnung bringt, so erhält man
für 190 Bögen folgende Ausbeute (III):
\begin{center}
\begin{tabular}{rrrr}
 & Sin. & Cos. & Tang. \\
0 & 89 & 126 & 87 \\
1 & 101 & 64 & 90 \\
2 & & & 12 \\
3 & & & 1
\end{tabular}
\end{center}
Will man diese Zahlen zu einer Abschätzung des Verhältnisses zwischen den
richtigen und unrichtigen Logarithmen anwenden, so muss man erst noch einen
kleinen Abzug machen. Es ist in (I) der Logarithm des Sinus von $45^{0}$ zweimal
gezählt, nemlich zugleich auch als $\log\cos45^{0}$; es ist ferner in der Tangentencolumne auch der Logarithm der Tangente von $45^{0}$ mitgezählt, der doch rational
ist. Dasselbe gilt von (III) und man muss daher, zu obigem Zweck, das Resultat davon so aussprechen, dass unter 568 geprüften irrationalen Logarithmen 301
sich als richtig, und 267 als unrichtig ausgewiesen haben. Dürfte man dies Verhältniss als durchschnittlich zutreffend betrachten, so würden unter den 68038
irrationalen Logarithmen des \textsc{Vega}'schen Thesaurus (indem man, wie billig, die
Cotangenten nicht mitzählt) nach der Wahrscheinlichkeit etwa 31983 fehlerhafte
anzunehmen sein.

Wahrscheinlich ist aber diese Zahl noch bedeutend zu klein. Ich finde in
meinen Papieren die vor längerer Zeit und zu andern Zwecken auf 14 Ziffern gemachte Berechnung der trigonometrischen Logarithmen für ein paar Gruppen von
Bögen, die nicht sprungsweise, sondern in denselben Intervallen wie die \textsc{Vega}'schen Tafeln, fortschreiten, woraus wenigstens hervorgeht, dass obiges Resultat
(III) noch keinen richtigen Maassstab für die Ungenauigkeit dieser Tafeln abgibt.
Während bei jenen 190 Bögen kein einziger Fall vorkommt, wo der Logarithm
eines Sinus oder eines Cosinus um mehr als \textit{eine} Einheit verbessert werden müsste,
sind solche Fälle gar nicht selten bei denjenigen Bögen, die \textit{nicht} in der Trigonometria Britannica vorkommen, und wo also das zur Vergleichung nöthige erst
durch besondere Rechnung herbeigeschafft werden muss. Ich füge daher die Resultate jener Rechnungen hier bei, da sie dazu dienen können, unserer Vorstellung von dem Grade der Ungenauigkeit der Zahlen in \textsc{Vega}'s Thesaurus eine festere Haltung zu geben.
